\documentclass[11pt,a4paper]{article}
\usepackage[T1]{fontenc}\usepackage[utf8]{inputenc}\usepackage{lmodern}
\usepackage[margin=2.3cm]{geometry}
\usepackage{amsmath,amssymb,bm}\usepackage{booktabs}
\usepackage[dvipsnames]{xcolor}\usepackage{mdframed}\usepackage{listings}
\usepackage[colorlinks=true,linkcolor=NavyBlue,urlcolor=NavyBlue]{hyperref}
\usepackage{microtype}

\definecolor{codebg}{HTML}{F4F4F0}\definecolor{rulec}{HTML}{C9C7BC}
\definecolor{keyc}{HTML}{7B3F00}\definecolor{warnc}{HTML}{8B2E1F}
\lstset{basicstyle=\ttfamily\scriptsize, columns=fullflexible, keepspaces=true,
  breaklines=true, showstringspaces=false, backgroundcolor=\color{codebg},
  xleftmargin=4pt, framexleftmargin=4pt, framextopmargin=3pt, framexbottommargin=3pt}
\newmdenv[linecolor=rulec,linewidth=0.6pt,backgroundcolor=codebg,innertopmargin=6pt,
  innerbottommargin=6pt,innerleftmargin=8pt,innerrightmargin=8pt,skipabove=8pt,skipbelow=8pt]{keybox}
\newmdenv[linecolor=warnc,linewidth=0.9pt,backgroundcolor=white,innertopmargin=6pt,
  innerbottommargin=6pt,innerleftmargin=8pt,innerrightmargin=8pt,skipabove=8pt,skipbelow=8pt]{caution}
\newcommand{\code}[1]{\texttt{\small #1}}
\newcommand{\vE}{\bm v_E}\newcommand{\nhat}{\hat{\bm n}}\newcommand{\that}{\hat{\bm t}}

\title{\vspace{-1.5cm}Drift consistency of the JOREK \code{model600} boundary conditions:\\
       an audit and an implementation plan}
\author{Máté Szűcs}
\date{27 August 2026}

\begin{document}
\maketitle

\begin{abstract}
\noindent
Switching on the weak Galerkin sheath $j$--$V$ boundary condition destroys the inner-leg
density enhancement, with shock-capturing diffusion excluded as the cause. This document
audits every boundary condition and fluid--kinetic coupling site in \code{model600} for
consistency with the $\bm E\times\bm B$ drift, against the SOLPS-ITER manual 3.0.9.

The result: of the five boundary conditions SOLPS requires to be used together for a
drift case, JOREK has one right, one half right, and three with no drift at all. The one
that is half right --- the Mach-1 condition --- contains a \emph{double unit conversion}
rather than a missing factor, and it is \emph{identically dormant} under the Dirichlet
$u$ boundary condition and \emph{activates} exactly when the sheath BC is switched on.
That is the same switch that produces the observed effect, which is the argument for
fixing it first. Section 7 gives the implementation plan; Section 8 the case for the
priority ordering; Section 9 the traps.
\end{abstract}

\section{The unit convention, confirmed five ways}

Everything below rests on one fact, so it is worth establishing beyond doubt:
\begin{equation}
  \boxed{\;V_{\rm par} = v_\parallel / |\bm B|\;}
\end{equation}
\begin{center}
\begin{tabular}{lll}
\toprule
location & expression & implies \\
\midrule
\code{mod\_boundary\_conditions.f90:617} & \code{direction/Btot * cs0} & physical speed $\div|B|$ \\
\code{mod\_expression.f90:1680} & \code{fact\_vpar = sqrt(BB2)/fact\_time} & $v_\parallel = V_{\rm par}|B|$ \\
\code{mod\_expression.f90:1487} & \code{Vsound = c\_s/sqrt(BB2)} & same, inverted \\
\code{mod\_particle\_wall\_interaction.f90:1770} & \code{vpar * norm2(B)} & explicit conversion \\
\code{mod\_fields.f90}, \code{calc\_vvector} & \code{v\_phi = F0*vpar/R} & consistent \\
\bottomrule
\end{tabular}
\end{center}

\section{The Mach-1 bug is a double conversion, not a missing factor}

The drift-compatible Bohm condition requires a correction $-\vE\!\cdot\!\nhat/b_n$ to
$v_\parallel$, which in $V_{\rm par}$ units is $-\vE\!\cdot\!\nhat/(b_n B_{\rm tot})$. The
geometric identity (derivation in \code{doc/jorek\_vpar\_bc\_drifts.pdf}) gives
\begin{equation}
  R^{2}\frac{\partial_b u}{\partial_b\psi}
  \;=\; -\frac{\vE\cdot\nhat}{b_n\,B_{\rm tot}} ,
  \label{eq:identity}
\end{equation}
with $\partial_b$ the \emph{tangential} derivative. The two sides are the same expression:
\textbf{$R^2\partial_b u/\partial_b\psi$ is already in $V_{\rm par}$ units}, the $B_{\rm tot}$
being generated by the $\partial_b\psi$ in the denominator. Line 617 divides it again:
\begin{lstlisting}
Mach1BC = - Vpar0 + direction/Btot * factor * cs0 + factor/Btot * BigR**2 * U0_b/ps0_b
!                                                         ^^^^^ converts an already-converted quantity
\end{lstlisting}

\begin{keybox}
The \code{cs0} term genuinely needs $/B_{\rm tot}$, because \code{cs0} is a physical speed.
The drift term does not, because it is not. Framing this as a \emph{double conversion}
rather than a \emph{missing factor} matters: the former is visible in the source and
verifiable against the identity, the latter invites the reasonable objection that the factor
might be absorbed elsewhere.
\end{keybox}

\noindent\textbf{Checked and excluded: it is not absorbed by the imposition.}
\code{boundary\_conditions\_add\_one\_entry} places $-\code{zbig}\cdot\code{Mach1BC\_v}$ on
the diagonal and $+\code{zbig}\cdot\code{Mach1BC}$ on the right-hand side, so \code{zbig}
cancels. It is imposed on \code{index(1)} (the value) and \code{index(iv\_dir)} (the
\emph{tangential} derivative). No perpendicular degree of freedom, no projection, no
mass-matrix normalisation --- nothing that could carry a field factor.

\section{SOLPS confirms the target form}

Manual 3.0.9, p.\ 413, \code{BCMOM=13}, titled \emph{``Drift-compatible, Bohm-Chodura sheath
entrance condition''}:

\begin{quote}\itshape
This condition imposes that the sum of the poloidal projection of the parallel velocity and
the poloidal $E\times B$ velocity (restricted to not exceed $\pm2c_{s,a}|b_x|$), equals at
least the specific species sound speed $c_{s,a}$, reprojected in the poloidal direction.
\end{quote}
\begin{equation}
  v_\parallel b_x + V^{(E\times B)}_{\rm pol} \;\ge\; c_{s,a} b_x
  \qquad\Longrightarrow\qquad
  v_\parallel \ge c_{s,a} - \frac{V^{(E\times B)}_{\rm pol}}{b_x}
\end{equation}
At a target the poloidal direction \emph{is} the wall normal, so $b_x\equiv b_n$ and
$V^{(E\times B)}_{\rm pol}\equiv\vE\!\cdot\!\nhat$. The $\bm E\times\bm B$ enters with
coefficient \textbf{1}: \code{b2mn.dat} exposes \code{'b2stbc\_cbc' '1.0'}, a multiplier on
the $E\times B$ velocity in the \code{BCMOM=13} case. There is no $1/B$ anywhere, and the
$\pm2c_s|b_x|$ clip is theirs.

\section{The audit}

\begin{center}
\begin{tabular}{llc}
\toprule
variable & boundary term & drift? \\
\midrule
\code{u} & sheath $j$--$V$ characteristic (\code{:605}) & \textbf{yes} (this \emph{is} the potential BC) \\
\code{vpar} & nodal \code{Mach1BC} (\code{mod\_boundary\_conditions:617}) & \textbf{half} (double-converted) \\
\code{rho} & \code{density\_reflection*r0*vpar0*ps0\_s} (\code{:661}) & \textbf{no} \\
\code{Ti,Te,T} & $(\gamma_{\rm sheath}\!-\!1)\cdot r_0 T\cdot$\code{vpar0}$\cdot$\code{ps0\_s} (\code{:666/670/673}) & \textbf{no} \\
\code{vpar} & integral-form Mach-1 (\code{:679}) & \textbf{no} (inactive by default) \\
\bottomrule
\end{tabular}
\end{center}

\noindent In SOLPS's terms: \code{BCPOT=11} correct, \code{BCMOM=13} half, \code{BCCON=14}
and \code{BCENE/BCENI=15} absent --- and the manual states these five must be used
\emph{in conjunction}.

\subsection{Why the missing fluxes are not a weak-form pairing error}

A boundary term is normally the surface term of a volume operator integrated by parts, so a
missing one can indicate a broken pairing. Here it does not. The $\bm E\times\bm B$
convection of density is in \emph{strong} form,
\begin{lstlisting}
mod_elt_matrix_fft.f90:1670   + v * BigR**2 * ( r0_s * u0_t - r0_t * u0_s ) * tstep * factor(var_rho,2)
\end{lstlisting}
--- test function times the bracket, \emph{not} its derivative. It was never integrated by
parts, produces no surface term, and none is missing for consistency. The code is internally
consistent; the drift boundary condition was simply never written.

\section{The particle module: charge exchange is right, collisions are not}

Two different background velocities are used for two processes in the same particle push.

\begin{center}
\begin{tabular}{lll}
\toprule
process & background velocity & status \\
\midrule
charge exchange (\code{:423}$\to$\code{:482}) & \code{calc\_vvector} $\to$ full $\vE + v_\parallel\hat{\bm b}$ & \textbf{correct} \\
Coulomb collisions (\code{:607}$\to$\code{:614}) & \code{P(1)*B/sim\%t\_norm} $\to$ parallel only & \textbf{bug} \\
\bottomrule
\end{tabular}
\end{center}

\noindent\code{calc\_vvector} builds \code{v\_R = -R*U\_Z + vpar*psi\_Z/R} etc., i.e.\ it
contains $\vE = R(-\partial_Z u,\partial_R u)$ explicitly.

\begin{keybox}
$|\vE|\approx170$\,m/s against $v_\parallel\sim c_s\sim3\times10^4$\,m/s looks like a
$0.6\%$ correction. It is not: $\vE\perp\bm B$ while $v_\parallel\parallel\bm B$, so omitting
it discards \textbf{100\% of the perpendicular background flow}. Kinetic impurity ions are
never collisionally entrained by the $\bm E\times\bm B$, and \code{mom\_par\_source} is
computed against the wrong reference frame.
\end{keybox}

\section{Fluid and kinetic recycling mirror each other}

\begin{center}
\begin{tabular}{ll}
\toprule
fluid, \code{mod\_boundary\_matrix\_open.f90:661} & \code{density\_reflection * r0 * vpar0 * ps0\_s} \\
kinetic, \code{mod\_particle\_wall\_interaction.f90:1770} & \code{Gamma\_d = n\_e*|vpar|*norm2(B)*cos\_alpha} \\
 & \code{\qquad\qquad + n\_e*c\_s*c\_angle} \\
\bottomrule
\end{tabular}
\end{center}
with \code{cos\_alpha} $=|\hat{\bm b}\!\cdot\!\nhat| = |b_n|$. Structurally identical, both
parallel-only.

\begin{caution}
\textbf{These must be changed together.} Fixing one alone breaks particle conservation across
the fluid--kinetic interface: the fluid would recycle a flux the kinetic side never creates,
or the reverse.
\end{caution}

\noindent The kinetic side is the easier of the two: the OpenMP private list at \code{:1738}
already carries \code{U} and \code{E}, and \code{calc\_RK4} returns \code{E}, so
$\vE = \bm E\times\bm B/B^2$ is directly available with no new field machinery.

\section{Implementation plan}

\begin{enumerate}\itemsep6pt

\item[\textbf{0.}] \textbf{Collision background flow $\to$ \code{vvector}.}
\code{particles/mod\_particle\_evolution.f90:614}. Replace
\code{P(1)*B/sim\%t\_norm} with \code{vvector}, already computed at \code{:423} in the same
iteration and already in m\,s$^{-1}$; the \code{interp\_PRZ([var\_Vpar])} at \code{:607}
then becomes redundant. Confirm \code{vvector} is computed unconditionally rather than only
inside the CX branch. \emph{Self-contained; no pairing. Makes collisions agree with charge
exchange, which is already correct.}

\item[\textbf{1.}] \textbf{\code{Mach1BC}: remove the spurious $/B_{\rm tot}$, add the clip.}
\code{models/model600/mod\_boundary\_conditions.f90}, behind
\code{bohm\_drift\_compatible = .false.}:
\begin{lstlisting}
:617   + factor * BigR**2 * U0_b/ps0_b               ! was factor/Btot * ...
:620   + factor * BigR**2 * element_size_0/ps0_b     ! Jacobian, same scaling
:655   + factor * BigR**2 * U0_bb/ps0_b              ! n_order >= 5 only
:656   + factor * BigR**2 * element_size_3/ps0_b     ! n_order >= 5 only
\end{lstlisting}
and clip the drift term to $\pm2\,\code{cs0}$. \emph{No pairing needed: the particle module
reads \code{vpar} through \code{calc\_NeTevpar} and inherits the corrected value. Leave
model401/502 untouched --- the flag defaults off.} Self-check: the \code{diag\_mach1} ratio
must go to $-\operatorname{sign}(b_n)$ when the flag is on.

\item[\textbf{2+3.}] \textbf{Recycling on the total flux, fluid \emph{and} kinetic together.}
Add $n\,(\vE\!\cdot\!\nhat)$ to \code{:661} and to \code{Gamma\_d} at \code{:1770}.

\item[\textbf{4.}] \textbf{Sheath heat transmission on the total flux.}
\code{:666}, \code{:670}, \code{:673}: apply $(\gamma_{\rm sheath}\!-\!1)$ to
$\Gamma_{\rm tot}$, not $\Gamma_\parallel$. Check the kinetic energy accounting
(\code{i\_wall\_heat\_in}) stays matched. \emph{This is the term that couples the density
asymmetry back into $T_e$, hence $\eta$, hence $E_\parallel$ --- the amplifier of the SOLPS
loop, which the boundary currently lacks entirely.}

\item[\textbf{5.}] \textbf{Diamagnetic velocity in the sheath condition.} SOLPS carries
$u^{(\rm dia)}_a b_z$ in \code{BCMOM=3}; JOREK carries it nowhere. Largest change, lowest
priority.
\end{enumerate}

\section{Why item 1 is the most pressing}

\begin{keybox}
Under \code{bcs(:)\%dirichlet\%u = .true.} (the default), $u$ is pinned along the wall, so
$\partial_b u = 0$ and therefore, by \eqref{eq:identity},
\[
   \text{the \code{Mach1BC} drift term is \textbf{identically zero}.}
\]
It contributes nothing, and the condition reduces to the plain Bohm condition
$v_\parallel = c_s$ --- which is correct, because with $\partial_b u = 0$ there is also no
$\bm E\times\bm B$ flux through the wall to correct for.

Release $u$ with the sheath BC and \emph{both} switch on together: the drift flux becomes
non-zero and the correction meant to compensate it activates --- \textbf{at half size}.

\textbf{The bug's activation switch is exactly the switch that produces the observed effect.}
\end{keybox}

\noindent And the error is asymmetric in the right direction to explain it. Writing the total
normal flux with the undersized correction,
\begin{equation}
  \bm v\cdot\nhat = v_\parallel b_n + \vE\cdot\nhat
  = c_s b_n + \underbrace{\Bigl(1-\tfrac{1}{B_{\rm tot}}\Bigr)}_{\approx\,0.5}\vE\cdot\nhat ,
\end{equation}
whereas the drift-compatible condition should give exactly $c_s b_n$. Half the drift flux
survives as excess, with the sign measured by the \code{diag\_mach1} diagnostic:

\begin{center}
\begin{tabular}{lccc}
\toprule
target & $\vE\!\cdot\!\nhat$ & correction should & result of halving it \\
\midrule
INNER & $+8.2\times10^{-7}$ (into wall) & \emph{reduce} $v_\parallel$ & too high $\to$ over-removes $\to$ \textbf{drain} \\
OUTER & $-1.4\times10^{-6}$ (out of wall) & \emph{raise} $v_\parallel$ & too low $\to$ under-removes $\to$ \textbf{accumulate} \\
\bottomrule
\end{tabular}
\end{center}

\noindent\textbf{It drains the high-field side and sources the low-field side} --- the exact
direction needed to remove an inner-leg density lobe, appearing exactly when the sheath BC
is enabled.

Items 2--3 push the \emph{same} way and therefore compound rather than cancel: recycling is
computed from $\Gamma_\parallel$, but under a correct drift-compatible condition
$v_\parallel b_n = c_s b_n - \vE\!\cdot\!\nhat$, so at the inner target ($\vE\!\cdot\!\nhat>0$)
the parallel part is smaller and the inner target is \emph{under}-recycled. Too much removed
and too little returned, both on the HFS.

\subsection*{Magnitudes, stated honestly}

In the mean the excess is $\sim0.14\%$ of the Bohm flux --- small. Near tangency the
diagnostic measured the required correction reaching $280\%$ of $c_s$, giving $\sim140\%$
excess locally, and it is resolvable on $24\%$ of the inner target's boundary weight against
$2.3\%$ of the outer's. So: small on average, large where it bites, and geometrically biased
toward the inner target.

\begin{caution}
This is a correct fix with the right sign in the right place, and it is \emph{not} a promise
that it explains the whole lobe. The dominant effect may still be volume redistribution by
the sheath-driven $\bm E\times\bm B$, which is real physics given that potential. What items
0--4 do is make the boundary consequences of the potential correct, so that the next
observation can be trusted.
\end{caution}

\section{Traps}

\begin{itemize}\itemsep3pt
\item \textbf{\code{c\_angle} may already be standing in for the drift.} It is a minimum-flux
floor at grazing incidence, and the drift correction is \emph{largest} at grazing incidence
because it goes as $1/b_n$. Adding the drift term without revisiting \code{c\_angle} risks
double counting precisely where both are biggest.
\item \textbf{\code{iv\_dir} orientation.} The fluid-side flux additions need
$\partial_b u$ with the correct \code{normal\_sign3} and \code{iv\_dir} sense. Getting it
wrong makes the term dissipative on some boundary elements and anti-dissipative on the rest
--- the same failure mode documented for \code{sheath\_wall\_vel} at
\code{mod\_boundary\_conditions.f90:1338}, which grows a mesh-scale alternating structure in
$\psi$, visible first in $zj$ and only later in $u$.
\item \textbf{Never change fluid recycling without the kinetic side}, or particle
conservation breaks across the interface.
\item \textbf{The clip is not optional.} The diagnostic measured a $0.89\%$ would-reverse
fraction; without the $\pm2c_s$ bound the corrected condition can reverse $v_\parallel$ near
tangency.
\end{itemize}

\vspace{1em}\hrule\vspace{0.6em}
\noindent\footnotesize
Branch \code{thermoelectric-ohm}, AUG \#38773, \code{model600}, $n_{\rm tor}=1$.
Reference: SOLPS-ITER manual 3.0.9 (\code{BCMOM}, pp.\ 407--423; \code{b2mn.dat} numerics).
Companions: \code{doc/jorek\_vpar\_bc\_drifts.pdf} (identity \eqref{eq:identity} and the
Mach-1 BC in full), \code{doc/sheath\_bc\_whitepaper.pdf} (the weak Galerkin $j$--$V$ BC),
\code{doc/hfshd\_findings.pdf} (campaign results).
\end{document}
